\section{The S-Chain Storage Market}
\label{sec:economics}

A decentralized storage substrate needs an economic model, both to price the
service to customers and to compensate the operators who physically store and
serve bytes. S-Chain needs two pricing models, because it serves two kinds of
data with opposite lifetimes, and a provider market underneath both.

\subsection{Two demand curves: metered and perpetual}

\paragraph{Metered (operational).} Hot, mutable, operational data --- a
Datastore tenant's active parts, a project's live SQLite file --- is priced
pay-as-you-go. This model exists today: the platform meters storage at a
per-gigabyte-month rate (with finer per-second and per-hour granularity) plus
egress, billed against a prepaid organization wallet, comparable to commodity
object-store pricing. It suits data whose value tracks its recency and whose
owner expects a recurring bill.

\paragraph{Perpetual (archival).} A second class of data should be stored
\emph{once} and kept \emph{forever} --- audit logs, published datasets, the
verifiable analytics commitments of Section~\ref{sec:chain}, content a user
wants to outlive their subscription. Recurring billing is the wrong instrument
here: it makes permanence contingent on an unbroken payment stream. The
Arweave-style alternative is an \emph{endowment}: the writer pays a single
upfront fee, a conservative fraction is reserved to cover immediate storage, and
the remainder capitalizes an endowment whose yield funds storage in perpetuity.
The model is solvent under one empirical assumption --- that the real cost of
storing a byte declines over time (the historical $\sim\!30\%$/year decline in
\$/GB) faster than the endowment is drawn down. Concretely, for a one-time fee
$F$, an annual real storage cost $c_0$ per byte declining at rate $r$, and an
endowment real return $\rho$, perpetual storage is funded when
\[
F \;\ge\; c_0 \sum_{t=0}^{\infty} \frac{(1-r)^t}{(1+\rho)^t}
\;=\; c_0\,\frac{1+\rho}{\rho + r},
\]
a finite sum precisely because $r>0$. Pricing $F$ at a deliberately pessimistic
$r$ (and re-pricing new writes as hardware costs fall) is what makes
``pay once, stored for two centuries'' actuarially sound rather than a promise.
This endowment does not exist yet and is the concrete economic build for S-Chain.

\subsection{The provider market underneath}

Both demand curves pay into a supply side: independent operators who run S-Chain
storage nodes (the Hanzo S3 / SeaweedFS volume servers) and earn for storing and
serving content-addressed objects. The accountability surface is already modeled
--- the network's service-node registry carries provider registrations, service
endpoints, and SLA/uptime records --- so the missing piece is the incentive
layer that turns those records into payments and penalties:

\begin{itemize}[leftmargin=1.4em]
  \item \textbf{Bonded providers.} An operator stakes a bond to join a storage
  pool, consistent with the network's bonded-signer model elsewhere; this gates
  a permissioned-but-open supply set rather than fully anonymous storage.
  \item \textbf{Proof of storage / availability.} Providers periodically prove
  they still hold the objects they are paid for (a challenge--response over
  content-addressed chunks), recorded as the registry's uptime/SLA proofs.
  \item \textbf{Payout and slashing.} The metered revenue and the endowment
  yield flow to providers in proportion to proven, served storage; failed
  availability proofs slash the bond. This is the same accounting that the
  platform's existing usage-metering and wallet rails perform for compute,
  extended to storage providers.
\end{itemize}

\subsection{A permissionless public good, paid in any currency}

The market above is deliberately \emph{not} operator-run or
us-in-control: S-Chain is intended as permissionless public-goods
infrastructure that web3 applications build on with no central leader, the way
one builds on a base layer rather than on a vendor. Two mechanisms make that
economically coherent.

\paragraph{Multi-currency fees routed through the DEX.} Applications should not
be forced to hold the network's base asset to pay for storage and compute. The
network's native DEX (the D-Chain) is the exchange layer: any L2/L3 lists its
coin into the DEX, and its applications pay infrastructure fees in \emph{their
own} currency, which is swapped natively --- through the relevant liquidity pool
(e.g.\ a ZOO/LUX pool) --- back into LUX to settle operator payouts. A storage
or compute provider is thus paid in the base asset while a tenant spends the
token it already holds; the DEX absorbs the currency conversion as a routine
swap. The fee economy is multi-currency at the edge and single-currency at
settlement.

\paragraph{Bootstrap inflation, then fees.} A storage market has a cold-start
problem: providers must be paid before fee volume exists. The resolution is a
\emph{temporary} inflation subsidy to seed provider rewards, tapering to zero as
fee revenue grows --- explicitly temporary because the base asset is
deflationary and cannot subsidize forever. Steady state is fee-funded: metered
leases and endowment yield, denominated in any listed currency and settled in
LUX, pay the bonded providers. Inflation is the on-ramp, not the engine.

\subsection{Why this closes the platform}

The economics complete the architecture rather than decorate it. The two query
engines of Section~\ref{sec:twoengine} both reduce, at the storage layer, to
content-addressed objects on S-Chain --- Datastore parts and encrypted SQLite
pages alike --- so a single market prices and endows \emph{all} platform data
uniformly: meter the hot working set, endow the permanent record, pay the
bonded providers who hold it, and let the same Quasar consensus that orders the
data's commitments also order the market's payments. A customer buys ``storage''
once; the platform decides per-object whether that is a recurring lease or a
perpetual endowment; and neither Datastore nor the decentralized SQL engine
needs its own billing system, because both are tenants of the one S-Chain
market.
